\chapter{Email Security}

\section{SMTP-Auth}
SMTP is the protocol used to send email. SMTP-Auth is an extension of ESMTP to provide authentication. It authenticates a client to the server before accepting messages from it, which is useful against spamming.

\textbf{Example 1: Authentication Failure}
\begin{enumerate}
    \item Start a connection request to the server \texttt{example.polito.it}.
    \item The client asks to use the extended version of the SMTP protocol by sending the message \texttt{EHLO}.
    \item The server responds with a positive message (code 250), indicating it speaks ESMTP and listing supported authentication methods.
    \item The client requests a method (e.g., PLAIN) that is \textbf{not} supported by the server.
    \item An error message is generated; authentication cannot be fulfilled.
\end{enumerate}

\textbf{Example 2: LOGIN Method Success}
\begin{enumerate}
    \item Start a connection request to the server \texttt{example.polito.it}.
    \item The client sends \texttt{EHLO}.
    \item The server responds (code 250) and lists supported methods: LOGIN, CRAM-MD5, DIGEST-MD5.
    \item The client chooses the LOGIN method and authenticates.
    \item \textbf{Note:} In LOGIN/PLAIN, the password is only encoded in BASE64. It is \textbf{not encrypted} and is visible if the channel is not protected by TLS.
\end{enumerate}

\section{ESMTP with TLS}
By default, data travels in cleartext between client and server. To add confidentiality, TLS can be used via the \texttt{STARTTLS} command. This upgrades a normal SMTP channel to a secure SMTP-over-TLS channel.

\begin{itemize}
    \item The client sends the \texttt{STARTTLS} command.
    \item A negotiation of security parameters (handshake) starts with the server.
    \item In case of success, the ESMTP transaction continues inside the encrypted tunnel.
    \item In case of failure (handshake fails or certificates are invalid):
    \begin{itemize}
        \item The client may quit (if configured to enforce security).
        \item Or, the connection may fall back to cleartext (opportunistic TLS).
    \end{itemize}
\end{itemize}

However, we cannot rely solely on channel security because we cannot guarantee that every node (hop) in the path is secure. We must protect the message itself (\textbf{End-to-End Security}).

Which kind of protection do we want?
\begin{itemize}
    \item \textbf{Integrity:} The message has not been modified in transit.
    \item \textbf{Authentication:} The receiver is able to identify the sender.
    \item \textbf{Non-repudiation:} The \textbf{sender} (not the server) cannot deny having sent a message.
    \item \textbf{Confidentiality:} The message should be readable only by the intended receiver.
\end{itemize}

\subsection{Secure MUA (Mail User Agent)}
The email system was originally designed only for text messages. We must ensure backward compatibility. The infrastructure (MTA, MSA, etc.) cannot be easily modified, so the \textbf{MUA} must be modified to handle security.

There are 4 main formats; they all require a secure MUA to verify or decrypt:
\begin{itemize}
    \item \textbf{Clear-signed:} The message body is in cleartext, but a digital signature is attached. Users without a secure MUA can still read the text (but cannot verify the signature).
    \item \textbf{Signed (Opaque):} The message and signature are encapsulated together. A secure MUA is required to view the message.
    \item \textbf{Encrypted (Enveloped):} The message is encrypted using a symmetric key, and that key is encrypted with the receiver's public key.
    \item \textbf{Signed and Enveloped:} A combination of the above; achieves all 4 security services.
\end{itemize}



\subsection{The Cryptographic Process}
\textbf{1. Canonicalization:}
Put the email message in a standard format (e.g., standardizing line endings to CR/LF) before hashing. This ensures the hash is consistent regardless of the computer system.

\textbf{2. Digital Signature (Integrity + Authentication + Non-Repudiation):}
\begin{itemize}
    \item Create a Hash (MIC - Message Integrity Code) of the message.
    \item Encrypt the Hash with the \textbf{Private Key of the Sender}.
    \item Send the message + the encrypted hash + the sender's certificate (containing the Public Key).
\end{itemize}

\textbf{3. Encryption (Confidentiality):}
\begin{itemize}
    \item A one-time \textbf{symmetric key} (Session Key) is generated.
    \item The message is encrypted with this symmetric key.
    \item The symmetric key is then encrypted with the \textbf{Public Key of the Receiver}.
    \item The package sent contains: The Encrypted Message + The Encrypted Session Key.
\end{itemize}

\section{Secure Email Standards}
\begin{itemize}
    \item \textbf{IETF Standards:} PEM (Obsolete), MOSS (Obsolete), S/MIME (The current industry standard).
    \item \textbf{Open Standards:} PGP (Pretty Good Privacy) and OpenPGP. Relies on a "Web of Trust" rather than centralized CAs.
    \item \textbf{Legacy/Gov:} X.400.
\end{itemize}

\textbf{Why MIME?}
RFC 822 (the original email standard) had limitations: designed only for English (7-bit ASCII) and had line-length limits.
MIME (Multipurpose Internet Mail Extensions) was introduced to allow:
\begin{itemize}
    \item Non-ASCII character sets.
    \item Binary attachments (images, video, etc.).
    \item Multipart bodies (recursiveness).
\end{itemize}
This ability to send binary data is essential for sending cryptographic keys and signatures.

\section{S/MIME (Secure/MIME)}
S/MIME is the dominant standard for commercial email security. It is architecturally based on:
\begin{itemize}
    \item \textbf{CMS (Cryptographic Message Syntax):} Based on PKCS \#7.
    \item \textbf{X.509v3:} For Certificates.
\end{itemize}

It supports algorithms for:
\begin{itemize}
    \item Hashing: SHA-256, SHA-512.
    \item Signature: RSA, ECDSA.
    \item Confidentiality: AES-128-CBC, AES-256-GCM.
\end{itemize}

\subsection{S/MIME Content Types}
New MIME types map email components to security functions:
\begin{itemize}
    \item \texttt{multipart/signed}: Used for \textbf{Clear-signed} messages. The body is clear; the signature is a separate attachment.
    \item \texttt{application/pkcs7-mime}: Used for \textbf{Enveloped} (encrypted) or \textbf{Opaque-signed} data. The content is a binary blob.
    \item \texttt{application/pkcs7-signature}: The format for the detached signature in a multipart/signed message.
\end{itemize}

\section{Mail Retrieval (POP/IMAP)}
Security is also required when retrieving mail from the server to the client.
\begin{itemize}
    \item \textbf{Authentication:} To ensure only the owner accesses the mailbox.
    \item \textbf{Confidentiality:} Basic POP3 and IMAP send passwords in cleartext.
\end{itemize}
To fix this, SSL/TLS wrappers are used (POP3S on port 995, IMAPS on port 993) or STARTTLS extensions within the standard ports.
